% ============================================================================
% LANGENTWURF - Sequenz 1a: Rally durchs Buch + Begrüßung
% ============================================================================
% TIER 1 (Vollständig) | Profil A (Gesamtschule heterogen)
% ============================================================================

\documentclass[11pt,a4paper]{article}

\usepackage[utf8]{inputenc}
\usepackage[T1]{fontenc}
\usepackage[ngerman]{babel}
\usepackage[left=2cm,right=2cm,top=2cm,bottom=2cm]{geometry}
\usepackage{helvet}
\usepackage{xcolor}
\usepackage{tabularx}
\usepackage{array}
\usepackage{enumitem}
\usepackage{fancyhdr}
\usepackage{lastpage}
\usepackage{tcolorbox}
\usepackage{booktabs}
\usepackage{amssymb}

\renewcommand{\familydefault}{\sfdefault}

\definecolor{spanisch}{HTML}{c41e3a}
\definecolor{grau}{HTML}{666666}
\definecolor{hellgrau}{HTML}{f0f0f0}
\definecolor{basis}{HTML}{4CAF50}
\definecolor{standard}{HTML}{2196F3}
\definecolor{erweiterung}{HTML}{9C27B0}

\pagestyle{fancy}
\fancyhf{}
\fancyhead[L]{\small\textcolor{grau}{Spanisch Klasse 7 -- Gesamtschule MV}}
\fancyhead[R]{\small\textcolor{grau}{LE-01a | Tier 1}}
\fancyfoot[C]{\small\thepage/\pageref{LastPage}}
\renewcommand{\headrulewidth}{0.5pt}

\newtcolorbox{kompetenzbox}{
    colback=hellgrau,
    colframe=spanisch,
    fonttitle=\bfseries,
    title=Kompetenzziele,
    arc=0pt,
    boxrule=1pt
}

\newtcolorbox{diffbox}{
    colback=white,
    colframe=grau,
    fonttitle=\bfseries,
    title=Differenzierung (nur Lehrkraft),
    arc=0pt,
    boxrule=0.5pt
}

\newtcolorbox{reflexbox}{
    colback=white,
    colframe=spanisch!50,
    fonttitle=\bfseries,
    title=Reflexion (nach der Stunde),
    arc=0pt,
    boxrule=0.5pt
}

\newcommand{\basis}{\textcolor{basis}{\textbf{[B]}}}
\newcommand{\standard}{\textcolor{standard}{\textbf{[S]}}}
\newcommand{\erweiterung}{\textcolor{erweiterung}{\textbf{[E]}}}

\begin{document}

% ============================================================================
% TITELSEITE
% ============================================================================
\begin{center}
{\Large\textbf{\textcolor{spanisch}{Langentwurf (Tier 1)}}}\\[0.5em]
{\LARGE\textbf{Rally durchs Buch + Begrüßung}}\\[0.3em]
{\large Sequenz 1a -- Empezamos, Std. 1--2}
\end{center}

\vspace{1em}

% ============================================================================
% KOPFBEREICH
% ============================================================================
\section*{Unterrichtsdaten}

\begin{tabularx}{\textwidth}{|l|X|l|X|}
\hline
\textbf{Fach} & Spanisch (2. FS) & \textbf{Datum} & \\
\hline
\textbf{Klasse} & 7 (Gesamtschule) & \textbf{Zeit} & 2 $\times$ 45 Min. \\
\hline
\textbf{Thema} & Orientierung + Begrüßung & \textbf{Raum} & \\
\hline
\textbf{Lehrwerk} & Qué pasa 1, Rally + S. 8--11 & \textbf{Profil} & A (heterogen) \\
\hline
\textbf{Sequenz} & 1 von 8 (Empezamos) & \textbf{SuS} & ca. 25 \\
\hline
\end{tabularx}

\vspace{1em}

% ============================================================================
% LERNVORAUSSETZUNGEN
% ============================================================================
\section*{Lernvoraussetzungen}

\begin{tabularx}{\textwidth}{|l|X|}
\hline
\textbf{Sprachlich} & Keine Vorkenntnisse in Spanisch; Englisch als 1. FS \\
\hline
\textbf{Methodisch} & Partnerarbeit bekannt; Umgang mit Lehrwerk vertraut (aus anderen Fächern) \\
\hline
\textbf{Inhaltlich} & Allgemeinwissen über Spanien/Lateinamerika variiert \\
\hline
\textbf{Heterogenität} & BR 20\%, MR 50\%, GY 30\% -- große Spannbreite \\
\hline
\end{tabularx}

\vspace{1em}

% ============================================================================
% KOMPETENZZIELE
% ============================================================================
\begin{kompetenzbox}

\textbf{Funktionale kommunikative Kompetenz:}
\begin{itemize}[nosep]
    \item \textbf{Hörverstehen:} Einfache Begrüßungsdialoge verstehen
    \item \textbf{Sprechen:} Begrüßungen und Verabschiedungen situationsgerecht anwenden
    \item \textbf{Leseverstehen:} Sich im Lehrwerk orientieren
\end{itemize}

\textbf{Sprachliche Mittel:}
\begin{itemize}[nosep]
    \item Redemittel: ¡Hola!, Buenos días, Buenas tardes, Buenas noches
    \item Redemittel: ¡Adiós!, ¡Hasta luego!, ¡Hasta mañana!
    \item Aussprache: Betonung, Vokale (ue in ``Buenos'')
\end{itemize}

\textbf{Interkulturelle kommunikative Kompetenz:}
\begin{itemize}[nosep]
    \item Spanischsprachige Länder auf der Weltkarte lokalisieren
    \item Begrüßungsformen in Spanien kennenlernen (Wangenkuss)
\end{itemize}

\textbf{Text- und Medienkompetenz:}
\begin{itemize}[nosep]
    \item Aufbau und Struktur des Lehrwerks verstehen
    \item Symbole und Verweise im Buch deuten
\end{itemize}

\textbf{Sprachlernkompetenz:}
\begin{itemize}[nosep]
    \item Internationalismen erkennen und nutzen
    \item Wörterverzeichnis als Hilfsmittel kennenlernen
\end{itemize}

\end{kompetenzbox}

\vspace{1em}

% ============================================================================
% STUNDENZIELE
% ============================================================================
\section*{Stundenziele}

\textbf{Stunde 1 -- Rally durchs Buch:}
\begin{enumerate}[nosep]
    \item Die SuS können sich im Lehrwerk orientieren (Aufbau, Symbole, Anhänge).
    \item Die SuS erkennen Internationalismen und erschließen deren Bedeutung.
    \item Die SuS entwickeln Neugier auf die spanische Sprache und Kultur.
\end{enumerate}

\textbf{Stunde 2 -- Begrüßung \& Verabschiedung:}
\begin{enumerate}[nosep]
    \item Die SuS können Begrüßungen situationsgerecht anwenden (Tageszeit).
    \item Die SuS können sich verabschieden (¡Adiós!, ¡Hasta luego!).
    \item Die SuS sprechen die Redemittel korrekt aus.
\end{enumerate}

\newpage

% ============================================================================
% VERLAUFSPLANUNG STUNDE 1
% ============================================================================
\section*{Verlaufsplanung Stunde 1: Rally durchs Buch}

\begin{tabularx}{\textwidth}{|c|c|X|c|c|l|}
\hline
\textbf{Phase} & \textbf{Zeit} & \textbf{Unterrichtsgeschehen} & \textbf{SF} & \textbf{Ziel} & \textbf{Material} \\
\hline
\multicolumn{6}{|l|}{\textbf{EINSTIEG}} \\
\hline
Begrüßung & 3' &
\begin{minipage}[t]{\linewidth}
\vspace{2pt}
L betritt Raum, begrüßt auf Spanisch:\\
``¡Hola! ¡Buenos días!''\\
SuS reagieren (vermutlich irritiert)
\vspace{2pt}
\end{minipage} & PL & Neugier & -- \\
\hline
Impuls & 5' &
\begin{minipage}[t]{\linewidth}
\vspace{2pt}
Weltkarte zeigen: ``Wo spricht man Spanisch?''\\
SuS sammeln Länder, L markiert
\vspace{2pt}
\end{minipage} & PL & IKK & Weltkarte \\
\hline
\multicolumn{6}{|l|}{\textbf{ERARBEITUNG}} \\
\hline
Hinführung & 5' &
\begin{minipage}[t]{\linewidth}
\vspace{2pt}
L zeigt Lehrwerk: ``Das ist euer Buch für Spanisch.''\\
Kurze Vorstellung der Lehrwerk-Geschichte (Ribadesella)
\vspace{2pt}
\end{minipage} & PL & Orient. & Buch \\
\hline
Rally & 20' &
\begin{minipage}[t]{\linewidth}
\vspace{2pt}
\textbf{Partnerarbeit:} Rally durchs Buch\\
-- Wie viele Unidades?\\
-- Wo ist das Wörterverzeichnis?\\
-- Welche Symbole gibt es?\\
-- Finde 5 spanische Wörter, die du verstehst\\
Wettbewerb-Charakter (wer zuerst fertig)
\vspace{2pt}
\end{minipage} & PA & Orient. & Rally-Blatt \\
\hline
\multicolumn{6}{|l|}{\textbf{SICHERUNG}} \\
\hline
Austausch & 7' &
\begin{minipage}[t]{\linewidth}
\vspace{2pt}
Ergebnisse zusammentragen\\
Internationalismen an Tafel sammeln:\\
hotel, restaurante, música, chocolate...
\vspace{2pt}
\end{minipage} & PL & Sichern & Tafel \\
\hline
Ausblick & 5' &
\begin{minipage}[t]{\linewidth}
\vspace{2pt}
L: ``Nächste Stunde lernen wir die ersten Wörter.''\\
Verabschiedung: ``¡Hasta luego!''
\vspace{2pt}
\end{minipage} & PL & Motivat. & -- \\
\hline
\end{tabularx}

\vspace{1em}

% ============================================================================
% VERLAUFSPLANUNG STUNDE 2
% ============================================================================
\section*{Verlaufsplanung Stunde 2: Begrüßung \& Verabschiedung}

\begin{tabularx}{\textwidth}{|c|c|X|c|c|l|}
\hline
\textbf{Phase} & \textbf{Zeit} & \textbf{Unterrichtsgeschehen} & \textbf{SF} & \textbf{Ziel} & \textbf{Material} \\
\hline
\multicolumn{6}{|l|}{\textbf{EINSTIEG}} \\
\hline
Begrüßung & 3' &
\begin{minipage}[t]{\linewidth}
\vspace{2pt}
L: ``¡Hola! ¡Buenos días!''\\
SuS wiederholen im Chor
\vspace{2pt}
\end{minipage} & PL & Aktivier. & -- \\
\hline
Wiederh. & 5' &
\begin{minipage}[t]{\linewidth}
\vspace{2pt}
Kurze Wiederholung: Internationalismen aus Std. 1\\
``Was haben wir letzte Stunde gelernt?''
\vspace{2pt}
\end{minipage} & PL & Anknüpf. & -- \\
\hline
\multicolumn{6}{|l|}{\textbf{ERARBEITUNG}} \\
\hline
Hörverstehen & 10' &
\begin{minipage}[t]{\linewidth}
\vspace{2pt}
Buch S. 10--11: Begrüßungsdialoge hören\\
SuS markieren gehörte Wörter\\
Tageszeiten zuordnen (Bilder: Sonne/Mond)
\vspace{2pt}
\end{minipage} & EA & HV & Audio, Buch \\
\hline
Systematisier. & 8' &
\begin{minipage}[t]{\linewidth}
\vspace{2pt}
Tafelanschrieb: Begrüßung vs. Verabschiedung\\
SuS übertragen in AB-01a Kasten 1
\vspace{2pt}
\end{minipage} & EA & Sichern & AB, Tafel \\
\hline
\multicolumn{6}{|l|}{\textbf{ÜBUNG}} \\
\hline
Chorsprechen & 5' &
\begin{minipage}[t]{\linewidth}
\vspace{2pt}
L spricht vor, Klasse wiederholt:\\
``¡Hola!'' -- ``¡Hola!''\\
``Buenos días'' -- ``Buenos días''...
\vspace{2pt}
\end{minipage} & PL & Ausspr. & -- \\
\hline
Mingle & 10' &
\begin{minipage}[t]{\linewidth}
\vspace{2pt}
\textbf{Mingle-Aktivität:}\\
SuS gehen durch den Raum\\
Bei Stopp (Musik aus): Nächsten begrüßen\\
L gibt Tageszeit vor (mañana/tarde/noche)
\vspace{2pt}
\end{minipage} & Klasse & Sprechen & Musik \\
\hline
\multicolumn{6}{|l|}{\textbf{SICHERUNG}} \\
\hline
AB-Arbeit & 7' &
\begin{minipage}[t]{\linewidth}
\vspace{2pt}
AB-01a Aufgabe 1: Zuordnung\\
Besprechung im Plenum
\vspace{2pt}
\end{minipage} & EA/PL & Sichern & AB \\
\hline
Abschluss & 2' &
\begin{minipage}[t]{\linewidth}
\vspace{2pt}
Verabschiedung: ``¡Adiós! ¡Hasta mañana!''\\
HA: AB Aufgabe 2
\vspace{2pt}
\end{minipage} & PL & Abschl. & -- \\
\hline
\end{tabularx}

\newpage

% ============================================================================
% DIFFERENZIERUNG
% ============================================================================
\section*{Differenzierung}

\begin{diffbox}

\textbf{Stunde 1 -- Rally:}

\begin{tabularx}{\textwidth}{|c|X|X|}
\hline
\textbf{Niveau} & \textbf{Maßnahme} & \textbf{Umsetzung} \\
\hline
\basis & Reduzierte Rally-Fragen (nur 4 statt 6) & Separates Blatt \\
\hline
\standard & Vollständige Rally & Standardblatt \\
\hline
\erweiterung & Zusatzaufgabe: 3 weitere Internationalismen finden & Mündlich \\
\hline
\end{tabularx}

\vspace{1em}

\textbf{Stunde 2 -- Begrüßung:}

\begin{tabularx}{\textwidth}{|c|X|X|}
\hline
\textbf{Niveau} & \textbf{Maßnahme} & \textbf{Umsetzung} \\
\hline
\basis & Nur ¡Hola! und ¡Adiós! produktiv & Weniger Varianten \\
\hline
\standard & Alle Begrüßungen produktiv & Volle Aufgabe \\
\hline
\erweiterung & Mini-Dialog frei formulieren & AB Aufgabe 4 \\
\hline
\end{tabularx}

\end{diffbox}

\vspace{1em}

% ============================================================================
% ERWARTETE SCHWIERIGKEITEN
% ============================================================================
\section*{Erwartete Schwierigkeiten}

\begin{tabularx}{\textwidth}{|l|X|X|}
\hline
\textbf{Bereich} & \textbf{Problem} & \textbf{Reaktion} \\
\hline
Aussprache & ``ue'' in ``Buenos'' als zwei Silben & Vorsprechen: [bwe-nos], Übertreibung \\
\hline
Grammatik & ``Buenos días'' vs. ``Buenas tardes'' (Genus) & Merksatz: ``der Tag'' = maskulin \\
\hline
Orthographie & Akzent bei ``días'' vergessen & Visuelle Hervorhebung an Tafel \\
\hline
Motivation & Schüchternheit bei Mingle & Erst PA, dann größere Gruppen \\
\hline
Tempo & Große Tempounterschiede bei Rally & Zusatzaufgaben für Schnelle \\
\hline
\end{tabularx}

\vspace{1em}

% ============================================================================
% TAFELANSCHRIEB
% ============================================================================
\section*{Tafelanschrieb}

\begin{center}
\fbox{\parbox{0.9\textwidth}{
\begin{center}
\textbf{\large Saludar y despedirse}\\
\textbf{Begrüßen und Verabschieden}
\end{center}

\vspace{0.5em}
\begin{tabular}{l|l}
\textbf{Begrüßung (Saludar)} & \textbf{Verabschiedung (Despedirse)} \\
\hline
¡Hola! (Hallo) & ¡Adiós! (Tschüss) \\
Buenos días (Guten Morgen) & ¡Hasta luego! (Bis später) \\
Buenas tardes (Guten Tag) & ¡Hasta mañana! (Bis morgen) \\
Buenas noches (Guten Abend) & \\
\end{tabular}

\vspace{0.5em}
\textit{Buenos} = maskulin (el día) | \textit{Buenas} = feminin (la tarde, la noche)
}}
\end{center}

\vspace{1em}

% ============================================================================
% MATERIALCHECKLISTE
% ============================================================================
\section*{Materialcheckliste}

\begin{tabularx}{\textwidth}{|l|X|c|c|}
\hline
\textbf{Material} & \textbf{Beschreibung} & \textbf{Std.} & \textbf{Anzahl} \\
\hline
Qué pasa 1 & Schülerbuch, Rally + S. 8--11 & 1--2 & Klassensatz \\
\hline
Rally-Blatt & Fragen zur Buchstruktur & 1 & Klassensatz \\
\hline
AB-01a & Arbeitsblatt Begrüßung & 2 & Klassensatz \\
\hline
Audio-CD & Track 1--3 (Begrüßungsdialoge) & 2 & 1 \\
\hline
Weltkarte & Spanischsprachige Länder & 1 & 1 \\
\hline
Musik & Für Mingle-Aktivität (beliebig) & 2 & 1 \\
\hline
\end{tabularx}

\vspace{1em}

% ============================================================================
% HAUSAUFGABE
% ============================================================================
\section*{Hausaufgabe}

\textbf{Nach Stunde 1:} Keine (erste Stunde)

\textbf{Nach Stunde 2:} AB-01a Aufgabe 2 (Tageszeiten zuordnen)

\vspace{1em}

% ============================================================================
% REFLEXION
% ============================================================================
\begin{reflexbox}

\textbf{Nach der Durchführung ausfüllen:}

\begin{enumerate}
    \item Wurde das Stundenziel erreicht?
    \begin{itemize}[nosep]
        \item[$\square$] vollständig
        \item[$\square$] größtenteils
        \item[$\square$] teilweise
        \item[$\square$] nicht erreicht
    \end{itemize}

    \item Wie war das Zeitmanagement?
    \begin{itemize}[nosep]
        \item[$\square$] passte genau
        \item[$\square$] zu wenig Zeit für: \rule{5cm}{0.4pt}
        \item[$\square$] zu viel Zeit für: \rule{5cm}{0.4pt}
    \end{itemize}

    \item Welche Schwierigkeiten traten auf?\\[2em]

    \item Was würde ich beim nächsten Mal anders machen?\\[2em]

    \item Differenzierung: War die Abstufung angemessen?
    \begin{itemize}[nosep]
        \item[$\square$] ja
        \item[$\square$] Basis zu schwer/leicht
        \item[$\square$] Erweiterung zu schwer/leicht
    \end{itemize}
\end{enumerate}

\end{reflexbox}

\vspace{1em}

% ============================================================================
% LERNSTANDSERHEBUNG
% ============================================================================
\section*{Lernstandserhebung}

\textbf{Beobachtungskriterien während des Unterrichts:}

\begin{tabularx}{\textwidth}{|X|c|c|c|}
\hline
\textbf{Kriterium} & \textbf{++} & \textbf{+} & \textbf{--} \\
\hline
SuS können Begrüßungen korrekt aussprechen & & & \\
\hline
SuS wählen situationsgerechte Begrüßung (Tageszeit) & & & \\
\hline
SuS beteiligen sich aktiv an Mingle-Aktivität & & & \\
\hline
SuS orientieren sich selbstständig im Lehrwerk & & & \\
\hline
\end{tabularx}

\vspace{0.5em}

\textbf{Minitest-Option (LP-01a oder MT-01a):} 10--15 BE, formativ

\vspace{1em}

% ============================================================================
% ANSCHLUSS
% ============================================================================
\section*{Anschluss: Stunde 3--4}

\textbf{Thema:} Sich vorstellen + Verb \textit{ser}

\textbf{Voraussetzung:} Begrüßungen sicher beherrscht

\textbf{Neuer Inhalt:}
\begin{itemize}[nosep]
    \item ¿Cómo te llamas? -- Me llamo...
    \item ¿De dónde eres? -- Soy de...
    \item Verb \textit{ser} (Singular)
\end{itemize}

\end{document}
